\chapter{Mengentheoretische Konstruktionen und Grenzen}
\label{b07guide:ch:constructions}

\section{Elementare Surjektionen}

\begin{theorem}[Identität und Inklusion]
Für jede Menge $A$ ist
$\operatorname{id}_A\colon A\twoheadrightarrow A$ surjektiv.
Für $A\subseteq B$ ist die Inklusion
$\iota\colon A\to B$, $a\mapsto a$, genau dann surjektiv,
wenn $A=B$.
\end{theorem}

\begin{theorem}[Projektionen und ihre leeren Randfälle]
\label{b07guide:thm:projections}
Für Mengen $A,B$ seien
\[
 \begin{aligned}
 \pi_1\colon A\times B\to A,&\qquad (a,b)\mapsto a,\\
 \pi_2\colon A\times B\to B,&\qquad (a,b)\mapsto b.
 \end{aligned}
\]
Dann gelten die genauen Kriterien
\[
 \begin{aligned}
 \pi_1\text{ surjektiv}&\Longleftrightarrow
                    A=\varnothing\ \text{ oder }\ B\ne\varnothing,\\
 \pi_2\text{ surjektiv}&\Longleftrightarrow
                    B=\varnothing\ \text{ oder }\ A\ne\varnothing.
 \end{aligned}
\]
\end{theorem}

\begin{theorem}[Schnitt mit einer festen Teilmenge]
Seien $C\subseteq A$. Dann ist
\[
 R_C\colon\mathcal P(A)\twoheadrightarrow\mathcal P(C),
 \qquad U\mapsto U\cap C
\]
surjektiv. Für $V\subseteq C$ ist seine Faser
\[
 R_C^{-1}[\{V\}]
   =\{V\cup W\mid W\subseteq A\setminus C\}.
\]
\end{theorem}

\section{Produkte und gemeinsame Zielwerte}

\begin{theorem}[Produkt zweier Surjektionen]
\label{b07guide:thm:product}
Seien $f\colon A\to B$ und $g\colon C\to D$. Die Funktion
\[
 f\times g\colon A\times C\to B\times D,\qquad
 (a,c)\mapsto(f(a),g(c))
\]
ist genau dann surjektiv, wenn
\[
 B\times D=\varnothing
 \quad\text{oder}\quad
 \bigl(f\text{ surjektiv und }g\text{ surjektiv}\bigr).
\]
Insbesondere ist das Produkt zweier Surjektionen surjektiv.
Bei nichtleerem $B\times D$ gilt auch die Umkehrung.
\end{theorem}

\begin{theorem}[Paarbildung mit gemeinsamem Definitionsbereich]
\label{b07guide:thm:pairing}
Für $f\colon A\to B$ und $g\colon A\to C$ ist die Funktion
\[
 H\colon A\to B\times C,\qquad a\mapsto(f(a),g(a))
\]
genau dann surjektiv, wenn
\[
 \forall b\in B\;\forall c\in C\quad
 f^{-1}[\{b\}]\cap g^{-1}[\{c\}]\ne\varnothing.
\]
Die Surjektivität beider Komponenten allein genügt im allgemeinen
nicht. Ist $H$ surjektiv und $B\times C\ne\varnothing$,
so sind beide Komponenten surjektiv.
\end{theorem}

\begin{theorem}[Surjektion nach einem Basiswechsel]
\label{b07guide:thm:basechange}
Seien $f\colon A\twoheadrightarrow B$ und $g\colon C\to B$.
Die Menge
\[
 P=\{(a,c)\in A\times C\mid f(a)=g(c)\}
\]
liefert die surjektive Funktion
\[
 p_C\colon P\twoheadrightarrow C,\qquad (a,c)\mapsto c.
\]
Ist auch $g$ surjektiv, so ist ebenfalls
$p_A\colon P\to A$, $(a,c)\mapsto a$, surjektiv.
Es wird keine zusätzliche Struktur auf $P$ vorausgesetzt.
\end{theorem}

\section{Markierte Vereinigungen}

\begin{definition}[Markierte Vereinigung]
Für eine gegebene Familie $(A_i)_{i\in I}$ sei
\[
 \coprod_{i\in I}A_i=\{(i,a)\mid i\in I,\ a\in A_i\}.
\]
Bei leerem $I$ ist diese Menge leer.
\end{definition}

\begin{theorem}[Markierte Vereinigung gegebener Surjektionen]
\label{b07guide:thm:taggedunion}
Sei $(f_i\colon A_i\to B_i)_{i\in I}$ eine gegebene Familie.
Dann ist
\[
 F\colon\coprod_{i\in I}A_i\longrightarrow\coprod_{i\in I}B_i,
 \qquad(i,a)\mapsto(i,f_i(a))
\]
genau dann surjektiv, wenn jedes $f_i$ surjektiv ist.
Diese Aussage gilt für beliebige Indexmengen, auch für
$I=\varnothing$. Die Funktionen $f_i$ sind bereits gegeben.
\end{theorem}

\begin{theorem}[Surjektion aus einer Mengenüberdeckung]
\label{b07guide:thm:forgettags}
Für jede gegebene Familie $(A_i)_{i\in I}$ ist
\[
 q\colon\coprod_{i\in I}A_i\twoheadrightarrow\bigcup_{i\in I}A_i,
 \qquad(i,a)\mapsto a
\]
surjektiv. Für eine Menge $B$ und eine gegebene Familie
$f_i\colon A_i\to B$ ist die Funktion
\[
 F\colon\coprod_{i\in I}A_i\to B,\qquad(i,a)\mapsto f_i(a)
\]
genau dann surjektiv, wenn
$\bigcup_{i\in I}f_i[A_i]=B$.
Eine Übereinstimmung der $f_i$ auf Überlappungen ist hier
nicht erforderlich.
\end{theorem}

\Needspace{30\baselineskip}
\section{Potenzmengen und eine Unmöglichkeitsaussage}

\begin{theorem}[Kanonische Urbilder für die Bildfunktion]
\label{b07guide:thm:powerset}
Sei $f\colon A\twoheadrightarrow B$. Dann ist
\[
 f_*\colon\mathcal P(A)\twoheadrightarrow\mathcal P(B),
 \qquad U\mapsto f[U]
\]
surjektiv. Mit
$f^*\colon\mathcal P(B)\to\mathcal P(A)$,
$V\mapsto f^{-1}[V]$, gilt
\[
 f_*\circ f^*=\operatorname{id}_{\mathcal P(B)}.
\]
Für jedes $V\subseteq B$ ist $f^{-1}[V]$ die größte
Teilmenge $U$ von $A$ mit $f[U]=V$.
Genauer gilt für $U\subseteq A$:
\[
 f[U]=V
 \quad\Longleftrightarrow\quad
 \left\{
 \begin{array}{l}
 U\subseteq f^{-1}[V],\\
 \forall b\in V\quad U\cap\operatorname{Fib}_f(b)\ne\varnothing.
 \end{array}\right.
\]
Eine kleinste solche Teilmenge wird nicht behauptet.
\end{theorem}

\begin{theorem}[Nichtleere Teilmengen und erneute Potenzmengenbildung]
\label{b07guide:thm:powersetiteration}
Jede Surjektion $f\colon A\twoheadrightarrow B$ liefert auch
die Surjektionen
\[
 \begin{aligned}
 \mathcal P(A)\setminus\{\varnothing\}
 &\twoheadrightarrow\mathcal P(B)\setminus\{\varnothing\},
 &U&\mapsto f[U],\\
 \mathcal P(\mathcal P(A))
 &\twoheadrightarrow\mathcal P(\mathcal P(B)),
 &\mathcal U&\mapsto\{f[U]\mid U\in\mathcal U\}.
 \end{aligned}
\]
Die erste Aussage gilt auch für $A=B=\varnothing$.
\end{theorem}

\begin{theorem}[Cantors Satz in Surjektivitätsform]
\label{b07guide:thm:cantor}
Für jede Menge $A$ gilt
\[
 \operatorname{Sur}(A,\mathcal P(A))=\varnothing.
\]
Es gibt also keine Funktion von $A$ auf sämtliche Teilmengen von $A$.
Allgemeiner gilt für $U\subseteq A$:
\[
 \operatorname{Sur}(U,\mathcal P(A))=\varnothing.
\]
Die Aussagen benötigen weder Endlichkeit noch eine
Mächtigkeitsordnung oder das Auswahlaxiom.
\end{theorem}